% Run GCT_KRON_20261120; GPT-6 Astra Pro; effort not exposed.
\section{Adaptive differential orders for thin partition shapes}

\subsection*{A-021. A linear degree gate for a bounded second row}
For $a\geq2$ and an integer $L\geq0$ define
\[
 B_a(L)=L+\sum_{j=2}^{a-1}\min(j,L).                    \tag{A-021.1}
\]
In particular $B_a(0)=0$; for $L\geq1$ the closed form is
\[
 B_a(L)=
 \begin{cases}
 aL-\binom L2-1,&1\leq L<a,\\
 L+\binom a2-1,&L\geq a.
 \end{cases}                                          \tag{A-021.2}
\]
For every finite-dimensional complex $V$, every $a\geq2$,
$b\geq1$, and $\lambda\vdash ab$, one has
\[
 b\geq B_a(\lambda_2)
 \quad\Longrightarrow\quad
 I(\operatorname{Ch}_b(V))_a[\lambda]=0.                \tag{A-021.3}
\]
As in A-020, missing second parts are zero and $[\lambda]$
means the entire isotypic component. Therefore for the actual
degree-$d$ determinant-closure multiplicity,
\[
 n\geq B_d(\lambda_2)
 \quad\Longrightarrow\quad
                K_n(\lambda)=a_\lambda(d[n])            \tag{A-021.4}
\]
for $d\geq2$. Since $B_d(L)\leq(d-1)L$, the simpler condition
$n\geq(d-1)\lambda_2$ is sufficient. These inequalities retain
all copies, and exclude an actual multiplicity obstruction for
every competing orbit closure in the same ambient space.

This refines A-020 by varying the exponent of the split binary
form. The adaptive choice was proposed by the root worker in
this run, and its proof was completed and checked here.
The two differential sequences originate in the cited 2026
Foulkes--Howe proof; no historical-priority claim is intended.

\begin{proof}
First handle one-row types. The degree-$a$ polynomial module
$\operatorname{Sym}^a((\operatorname{Sym}^b V)^*)$ contains
exactly one copy of $S_{(ab)}V^*$: the top-weight monomial is
the $a$-th power of the coefficient of the pure $b$-th power.
Restriction to pure powers $v^b$ is nonzero on this copy, since
that monomial evaluates to $\ell(v)^{ab}$ for a nonzero linear
form $\ell$. Equivariance and irreducibility imply that the
one-row copy has zero intersection with the Chow ideal. This
proves (A-021.3) for $\lambda_2=0$, in every dimension where
the type occurs.

Assume henceforth $L=\lambda_2\geq1$. The functions $B_a(L)$
are nondecreasing in $L$, satisfy $B_2(L)=L$, and obey
\[
 B_a(L)=B_{a-1}(L)+\min(a-1,L)\quad(a\geq3),
 \qquad B_a(L)\geq L.                                 \tag{A-021.5}
\]
We induct on $a$. At $a=2$ there are no quadratic Chow
equations for any $b\geq1$: for $b\geq2$ this is the cited
Foulkes--Howe Proposition 3, and for $b=1$ the Chow cone is
the whole vector space. Also any partition of $2b$ satisfies
$\lambda_2\leq b$, consistently with the claimed base bound.

Let $a\geq3$, put
\[
 r=a-1,\qquad s=\min(r,L),\qquad m=b-s.
\]
Then $s\geq1$ and the hypothesis gives
\[
                   m\geq B_r(L)\geq L\geq1.            \tag{A-021.6}
\]
Take $T$ in the full $\lambda$ isotypic component of the
symmetric $a$-linear forms polarizing the Chow ideal.
For $Q=y_1\cdots y_m$, binary factorization of
$x^s+ct^s$ and extraction of the coefficient of $c$ give
\[
                     T(x^sQ,\ldots,x^sQ,t^sQ)=0.       \tag{A-021.7}
\]
Keep all $x,t,y_i$ as independent vector-variable blocks, and
set $H=T(x^sQ,\ldots,x^sQ,t^b)$. Apply successively
$D_{y_i,t}=\sum_j(y_i)_j\partial/\partial t_j$ for
$i=1,\ldots,m$. Their composition is $b!/s!$ times
(A-021.7), so is zero. Before the $i$-th derivative, the
source $t$ has degree $p=b-i+1\geq s+1$, whereas the new
destination $y_i$ has degree $q=r$.

If $s=r$, then $p>q$, giving injectivity by the differential
kernel lemma (A-020.3). If $s=L$, then $p>L$; the same lemma
gives injectivity on type $\lambda$ even if $p\leq q$.
These cases cover $s=\min(r,L)$, including equality. Each
intermediate polynomial has only type $\lambda$, since both
evaluation and differentiation are equivariant. Thus every
derivative is injective on the image under consideration, and
$H=0$.

The equivariant map
\[
 T\longmapsto\bigl[(x,t,Q_1,\ldots,Q_r)\longmapsto
                         T(x^sQ_1,\ldots,x^sQ_r,t^b)\bigr]
\]
therefore has image in
\[
 \operatorname{Sym}^{rs}V^*\otimes\operatorname{Sym}^bV^*
                         \otimes I(\operatorname{Ch}_m(V))_r.
                                                               \tag{A-021.8}
\]
Indeed its diagonal in the $Q_i$ vanishes on every product,
and expansion in the independent $x,t$ coordinates puts every
coefficient in that homogeneous Chow ideal. By induction,
(A-021.6), and monotonicity of $B_r$, the last factor contains
no type $\nu$ with $\nu_2\leq L$: for $\nu_2>0$ one has
$m\geq B_r(L)\geq B_r(\nu_2)$, while the one-row case
was already proved. Two horizontal Pieri steps for the two
symmetric-power factors only add boxes. Hence (A-021.8) has
no type $\lambda$, whose second row has length $L$.
Equivariance forces the image of our $T$ to be zero. We have
obtained
\[
 T(x^sQ_1,\ldots,x^sQ_r,t^b)=0
                  \quad(Q_i\in\operatorname{Sym}^mV). \tag{A-021.9}
\]

Now put $f=T(v_1^b,\ldots,v_r^b,t^b)$ and introduce the
independent block $x$. Applying $D_{x,v_i}^s$ successively
for $i=1,\ldots,r$ gives
\[
 \left(\frac{b!}{m!}\right)^r
                  T(x^sv_1^m,\ldots,x^sv_r^m,t^b)=0
\]
by (A-021.9). Before any individual derivative the source
$v_i$ has degree at least $b-s+1=m+1>L$ by (A-021.6).
Thus (A-020.3) makes each derivative injective on type
$\lambda$, regardless of the growing degree in $x$.
The composition is injective on the image of $T$, so $f=0$.
As pure powers span every input space, this implies $T=0$.
The induction proves (A-021.3).

The closed form (A-021.2) is the sum in (A-021.1): when
$L<a$, sum $1+\cdots+L$ followed by $a-1-L$ copies of $L$,
remove the $j=1$ term and then add the leading $L$.
When $L\geq a$, all $j$ in the sum are at most $L$.
The bound $B_a(L)\leq(a-1)L$ follows because its $a-1$
terms are each at most $L$ (and also holds for $L=0$).
Finally, Chow products are diagonal determinants, so the
ideal containment and semisimple quotient argument in A-020
give (A-021.4), retaining every ambient multiplicity copy.
\end{proof}

\paragraph{Meaning of the strengthened gate.}
For $1\leq\lambda_2=L<d$, the sufficient bound is
\[
                          n\geq dL-\binom L2-1.
\]
For a fixed second row this grows linearly in $d$, whereas
the uniform Foulkes--Howe condition grows quadratically.
Combining the two verified bounds, an actual multiplicity
obstruction must satisfy
\[
        d(d-1)>n,\qquad B_d(\lambda_2)>n,
        \quad\hbox{and hence}\quad (d-1)\lambda_2>n.    \tag{A-021.10}
\]
All previous padding, row-number, and ambient-multiplicity
conditions remain applicable. This is a restriction on actual
orbit-closure multiplicities, without asserting support
saturation for Kronecker coefficients.

The scheme and comodule interpretation is exactly A-020's:
the maps in (A-021.8) are equivariant maps of Schur and
symmetric functors on locally free modules, the differential
operators are integral, and the proof is carried out after
base change to $\mathbb C$. The exact Grothendieck locators
are SGA 3, I, Propositions 4.7.2--4.7.3, pp.~34--36, and
SGA 6, VI, Theorem 3.3, pp.~392--393. The use of polynomial
spectators, rather than fixed choices of $x,t$, is what
preserves the full comodule type throughout the argument.
